\documentclass[11pt,a4paper]{article}
\usepackage[margin=2.3cm]{geometry}
\usepackage{booktabs,tabularx,enumitem,listings,xcolor}
\usepackage[hidelinks]{hyperref}

\lstset{basicstyle=\ttfamily\small,backgroundcolor=\color{gray!8},frame=single,
        rulecolor=\color{gray!40},breaklines=true,columns=fullflexible}
\setlist{nosep}
\newcommand{\svc}[1]{\texttt{#1}}

\title{Laboratory 1: Short project\\[2pt]\large Docker, Compose and Swarm}
\author{Marc De Rialp \and David Franco \and Omar Antonio Cornejo \and David Madrid \and Victor Yakovlev\\[4pt]
\small Cloud Computing and Distributed Systems\\ MSc in Machine Learning and Cybersecurity for Internet Connected Systems, UPC}
\date{October 2026}

\sloppy
\begin{document}
\maketitle

\section{Introduction}

We containerised the infrastructure of Figure~1 of the statement: the shared services
(\emph{basic}), a company application (\emph{app}), an observability platform
(\emph{stats}) and Moodle (\emph{atenea}). As the core application we wrote a small review
site in Flask: people leave a comment with a rating from 1 to 5, the average is cached in
Redis, each new comment sends a mail through our SMTP relay and an automatic feeder
collects ratings published in other web pages.

The repository contains \texttt{compose.yml} (local), \texttt{production.yml} (swarm),
one folder per image with its Dockerfile and files, and this document.

\section{Architecture}

Table~\ref{tab:services} maps the boxes of the diagram to our services. Each stack is an
overlay network, and a service only joins the networks it needs. For example the
databases are not in the \texttt{proxy} network, so the frontend can not reach them.

\begin{table}[h]
\centering\small
\begin{tabularx}{\textwidth}{llX}
\toprule
Stack & Service & Description \\
\midrule
basic  & \svc{frontend}    & nginx, the only public service (80/443), proxies by subdomain \\
       & \svc{letsencrypt} & certbot, gets and renews one certificate for all subdomains \\
       & \svc{registry}    & Docker registry with htpasswd authentication \\
       & NFS               & NFSv4 server (port 2049), shared disk for the nodes \\
\midrule
app    & \svc{core}        & Flask + gunicorn application \\
       & \svc{static}      & Apache serving \texttt{/static/} \\
       & \svc{db}          & PostgreSQL 16 (databases \texttt{app} and \texttt{grafana}) \\
       & \svc{redis}       & cache of the average rating \\
       & \svc{smtp}        & Postfix relaying the mails to an external provider \\
       & \svc{auto}        & feeder: reads schema.org ratings of web pages \\
       & \svc{backup}      & cron that dumps the databases and moodledata \\
\midrule
atenea & \svc{moodle}, \svc{moodle-cron} & Moodle 4.5 LTS (web and cron, same image) \\
       & \svc{moodle-db}   & MariaDB 11.4 \\
\midrule
stats  & \svc{grafana}     & dashboards, with the Infinity plugin (the ``?'') \\
       & \svc{prometheus}, \svc{cadvisor} & metrics of the containers and the core app \\
\bottomrule
\end{tabularx}
\caption{Services of the infrastructure.}
\label{tab:services}
\end{table}

Public names in production: \texttt{app.<domain>}, \texttt{lms.<domain>},
\texttt{grafana.<domain>} and \texttt{registry.<domain>}. Locally the same names are used
with \texttt{localhost} on port 8080 (browsers resolve \texttt{*.localhost} to 127.0.0.1).

\section{Images and design decisions}

\paragraph{Frontend.} Based on \texttt{nginx:alpine}. We use the template feature of the
official image (\texttt{envsubst} on startup) with two templates, local (http) and
production (https), selected with \texttt{NGINX\_ENVSUBST\_TEMPLATE\_DIR}. Upstreams are
resolved at request time through the Docker DNS, so nginx starts even if a backend is not
running yet. nginx refuses to start without a certificate, so a small script creates a self
signed one and switches to the Let's Encrypt one as soon as certbot gets it.

\paragraph{Let's Encrypt.} The official certbot image with a loop: \texttt{certbot certonly
--webroot --keep-until-expiring} every 12 hours, which issues the certificate the first time
and renews it later. The challenge files are written in a volume shared with nginx.

\paragraph{Core app.} Python 3.12 slim, runs as a non-root user. It creates its tables on
the first request, so there is no separate migration step. It exposes \texttt{/metrics} for
Prometheus (blocked from outside in nginx) and \texttt{/api/stats} with JSON for Grafana.
The page includes the average rating as JSON-LD (schema.org), the same format online shops
use for search engines.

\paragraph{Auto.} The feeder downloads the pages in \texttt{FEED\_URLS}, extracts the
\texttt{aggregateRating} of the JSON-LD and posts it to \texttt{/api/ratings}, protected by
a bearer token. By default it reads our own app, but any page with that markup works.

\paragraph{SMTP.} Postfix on Alpine. It only accepts mail from the Docker networks and
relays it to the provider configured in \texttt{.env} with SASL authentication and
mandatory TLS. In the local environment it relays to mailpit, so the mails can be checked in
the browser.

\paragraph{Moodle.} \texttt{php:8.3-apache} with the extensions Moodle needs and the 4.5 LTS
release downloaded at build time. \texttt{config.php} reads the URL from the environment and
the password from a Docker secret. The entrypoint waits for MariaDB and installs Moodle from
the CLI the first time; to avoid two replicas installing at the same time, the one that
manages to create a lock directory in the shared moodledata does it and the rest wait.
\svc{moodle-cron} is the same image running \texttt{admin/cli/cron.php} every minute.

\paragraph{Grafana.} Configured as code: the datasources, the dashboard provider and the
dashboard (JSON) are copied into the image, so the whole configuration is versioned in git.
For the ``?'' we chose the \textbf{Infinity} datasource plugin, which turns any JSON, CSV or
XML endpoint into a datasource. We use it to show the average rating and the ratings
captured by the feeder directly from our API, without writing an exporter.

\section{Deployment}

\subsection{Cluster}

We used three Ubuntu 24.04 virtual machines in the same private network, all of them swarm
managers so the cluster keeps working if one node fails (quorum of 2 out of 3). Node~1 has
the label \texttt{db=true} and also runs the NFS server. A wildcard DNS record
\texttt{*.<domain>} points to the nodes; thanks to the routing mesh any node accepts the
traffic and forwards it to a frontend replica.

\begin{lstlisting}
docker swarm init --advertise-addr 10.0.0.4      # node1, then join node2 and node3
docker node update --label-add db=true node1
\end{lstlisting}

\subsection{Shared storage}

The NFS server needs a privileged container and swarm services can not be privileged, so it
is the only service started with plain Compose on the storage node
(\texttt{basic/nfs/nfs.yml}). It exports \texttt{/srv/nfs} with NFSv4 only, restricted to
the cluster subnet. In \texttt{production.yml} the shared volumes use the \texttt{local}
driver with \texttt{type: nfs}, so Docker mounts them automatically on whatever node a task
is scheduled. Nothing is mounted by hand on the hosts.

\subsection{Compose files}

\begin{itemize}
\item \texttt{compose.yml}: local environment, no TLS, named volumes only, mailpit instead
of a real mail provider.
\item \texttt{production.yml}: the whole infrastructure for the swarm. It keeps the
\texttt{build:} keys, so the same file is used to build and push the images
(\texttt{docker stack deploy} ignores them).
\end{itemize}

\subsection{Steps}

\begin{lstlisting}
set -a; . ./.env; set +a                # docker stack does not read .env
docker compose -f production.yml build
docker compose -f production.yml push
docker stack deploy --with-registry-auth -c production.yml lab1
\end{lstlisting}

The images are pulled from our registry at \texttt{127.0.0.1:5000}. This address works on
every node because the registry port is published on the routing mesh, and Docker accepts
plain HTTP for localhost. On the very first deployment the registry does not exist yet, so we
deploy once (the registry starts and the other services wait), push the images and deploy
again. \texttt{--with-registry-auth} forwards the registry login to the other nodes.

\section{Scalability}

No service uses \texttt{container\_name}, fixed IPs or host ports other than the routing
mesh, so all of them can run with several containers. Table~\ref{tab:scale} shows the
replicas we use and why some are kept at one.

\begin{table}[h]
\centering\small
\begin{tabularx}{\textwidth}{llX}
\toprule
Service & Replicas & Reason \\
\midrule
\svc{core} & 3 & stateless, data in PostgreSQL and Redis \\
\svc{frontend}, \svc{static}, \svc{smtp} & 2 & stateless \\
\svc{registry} & 2 & storage on NFS, same \texttt{REGISTRY\_HTTP\_SECRET} \\
\svc{moodle} & 2 & moodledata on NFS, sessions in the database \\
\svc{grafana} & 2 & its state is in PostgreSQL, not in a local SQLite \\
\svc{cadvisor} & global & one per node \\
\svc{db}, \svc{moodle-db} & 1 & DBMS without automatic scaling \\
\svc{redis}, \svc{prometheus} & 1 & in-memory cache / local time series database \\
\svc{letsencrypt}, \svc{backup}, \svc{moodle-cron}, \svc{auto} & 1 & periodic jobs \\
\bottomrule
\end{tabularx}
\caption{Replicas per service.}
\label{tab:scale}
\end{table}

Updates use \texttt{order: start-first} and \texttt{failure\_action: rollback}, so a new
version only replaces the old one if it starts correctly. Prometheus discovers the replicas
of the core app with the DNS name \texttt{tasks.core}, so after
\texttt{docker service scale lab1\_core=5} the new containers appear in Grafana without
changes. The footer of the app shows the container that answered, which makes the load
balancing easy to see.

\section{Data management and backups}

\begin{table}[h]
\centering\small
\begin{tabular}{lll}
\toprule
Data & Volume & Location \\
\midrule
PostgreSQL, MariaDB, Prometheus & \texttt{db}, \texttt{moodle-db}, \texttt{prometheus} & local disk of node 1 \\
moodledata & \texttt{moodledata} & NFS \\
certificates, ACME challenges & \texttt{certs}, \texttt{acme} & NFS \\
registry images & \texttt{registry} & NFS \\
backups & \texttt{backups} & NFS \\
Redis & none & memory only (cache) \\
\bottomrule
\end{tabular}
\caption{Where the data is stored.}
\end{table}

We decided not to put the databases on NFS: database engines on network filesystems risk
corruption with locking problems or network cuts, and they are slow. They are pinned with a
placement constraint to the node labelled \texttt{db=true} and use local volumes; what goes
to the shared disk are the dumps.

The \svc{backup} service is an Alpine image with \texttt{crond}. Every night at 03:30 it
creates a folder in the backup volume with:
\begin{itemize}
\item \texttt{app.dump} and \texttt{grafana.dump}: \texttt{pg\_dump -Fc} (custom format,
allows restoring single tables);
\item \texttt{moodle.sql.gz}: \texttt{mariadb-dump --single-transaction}, consistent
without locking the tables;
\item \texttt{moodledata.tar.gz}: moodledata without caches, sessions and temporary files.
\end{itemize}
Folders older than 14 days are deleted. A backup can be launched by hand with
\texttt{docker exec <backup container> /backup.sh}. Restoring, from node 1:

\begin{lstlisting}
docker exec -i <db container> pg_restore -U app -d app --clean < /srv/nfs/backups/<date>/app.dump
gunzip -c /srv/nfs/backups/<date>/moodle.sql.gz | docker exec -i <moodle-db container> \
  sh -c 'mariadb -u moodle -p"$(cat /run/secrets/moodle_db_password)" moodle'
tar -xzf /srv/nfs/backups/<date>/moodledata.tar.gz -C /srv/nfs/moodledata
\end{lstlisting}

The main limitation is that the backups are on the same server as the shared data. For a
real deployment we would copy them off-site (for example with restic to object storage) and
encrypt them. Also, if node~1 fails, the databases stop until they are restored on another
node from the last dump.

\section{Security}

\begin{itemize}
\item Only nginx is exposed; port 80 only serves ACME challenges and redirects to HTTPS
(TLS 1.2/1.3 and HSTS).
\item Passwords are Docker secrets mounted in \texttt{/run/secrets}, never written in the
compose files; the \texttt{secrets/} folder is ignored by git.
\item The overlay networks \texttt{proxy}, \texttt{app} and \texttt{atenea} are encrypted
(IPsec between nodes), and each stack is isolated in its own network.
\item The SMTP relay is not published and only accepts mail from the internal networks.
\item The registry requires authentication (bcrypt htpasswd) and \texttt{/metrics} is not
reachable from outside.
\end{itemize}
Known weak points: ports published by Docker skip the host firewall (ufw), so the registry
port 5000 must be closed in the cloud provider firewall; and the NFS export uses
\texttt{no\_root\_squash} because some containers write as root, which we limited to the
cluster subnet.

\section{Conclusions}

The whole infrastructure is deployed with one command per step (build, push, deploy) and
every service can run on any node, except the databases, which we pinned on purpose. The
part that took the most time was Moodle (installation with several replicas and the shared
moodledata) and making nginx start before the real certificate existed.

\subsection*{Work distribution}
\begin{tabularx}{\textwidth}{lX}
Member 1 & basic stack: frontend, Let's Encrypt, registry, NFS \\
Member 2 & core app, static files, Redis \\
Member 3 & databases, SMTP relay, feeder, backups \\
Member 4 & Atenea (Moodle and MariaDB) \\
Member 5 & stats stack and swarm cluster \\
\end{tabularx}

\end{document}
